\documentclass[12pt]{article}

% ============================================================
% PACKAGES
% ============================================================

\usepackage{amsmath}      % Mathematical environments: align, equation, etc.
\usepackage{amsthm}       % Theorem, definition, proof environments
\usepackage{amsfonts}     % Mathematical fonts such as \mathbb
\usepackage{amssymb}      % Additional mathematical symbols
\usepackage{mathtools}    % Extensions to amsmath
\usepackage{geometry}     % Page margins
\usepackage{enumitem}     % Customize enumerate/itemize
\usepackage{mdframed}     % Framed boxes
\usepackage{tikz}         % Draw diagrams
\usepackage{hyperref}     % Clickable references and links

\geometry{
    letterpaper,
    margin=0.8in
}

\linespread{1.2}


% ============================================================
% NUMBER SYSTEMS
% ============================================================
% Instead of writing \mathbb{R} every time, use \RR.

\newcommand{\RR}{\mathbb{R}}   % Real numbers
\newcommand{\CC}{\mathbb{C}}   % Complex numbers
\newcommand{\NN}{\mathbb{N}}   % Natural numbers
\newcommand{\ZZ}{\mathbb{Z}}   % Integers
\newcommand{\QQ}{\mathbb{Q}}   % Rational numbers


% ============================================================
% COMMON MATHEMATICAL NOTATION
% ============================================================
% These commands are useful for frequently used notation.

\newcommand{\set}[1]{\left\{#1\right\}}
% Example:
%   \set{x \in \RR : x > 0}
% produces:
%   { x in R : x > 0 }

\newcommand{\abs}[1]{\left|#1\right|}
% Example:
%   \abs{x-y}
% produces:
%   |x-y|

\newcommand{\norm}[1]{\left\|#1\right\|}
% Example:
%   \norm{x}
% produces:
%   ||x||

\newcommand{\inner}[2]{\left\langle #1,#2\right\rangle}
% Example:
%   \inner{x}{y}
% produces:
%   <x,y>

\newcommand{\parens}[1]{\left(#1\right)}
% Example:
%   \parens{x+y}

\newcommand{\brac}[1]{\left[#1\right]}
% Example:
%   \brac{a,b}

\newcommand{\angles}[1]{\left\langle #1\right\rangle}
% Example:
%   \angles{x,y}


% ============================================================
% THEOREM / THEORY COMPONENTS
% ============================================================

% "definition" style is upright text.
% Use this for definitions, examples, and remarks.

\theoremstyle{definition}

% ------------------------------------------------------------
% DEFINITION
% ------------------------------------------------------------
% Use when introducing a formal mathematical definition.
%
% Usage:
%
% \begin{definition}[Metric]
% A metric on a set X is a function ...
% \end{definition}
%
% [Metric] is optional and gives the definition a name.

\newtheorem{definition}{Definition}[section]


% ------------------------------------------------------------
% EXAMPLE
% ------------------------------------------------------------
% Use to give a concrete example of a definition/theorem.
%
% Usage:
%
% \begin{example}
% Consider X = \RR with the usual distance.
% \end{example}

\newtheorem{example}[definition]{Example}


% ------------------------------------------------------------
% REMARK
% ------------------------------------------------------------
% Use for observations, comments, or small additional facts.
% It is NOT intended for major mathematical results.
%
% Usage:
%
% \begin{remark}
% This definition does not require X to be a subset of \RR^n.
% \end{remark}

\newtheorem{remark}[definition]{Remark}


% ============================================================
% THEOREM-LIKE COMPONENTS
% ============================================================

% "plain" style makes the mathematical statements italic.

\theoremstyle{plain}


% ------------------------------------------------------------
% THEOREM
% ------------------------------------------------------------
% Use for an important/general mathematical result.
%
% Usage:
%
% \begin{theorem}[Cauchy--Schwarz Inequality]
% For all x,y \in \RR^n,
% \[
% |\inner{x}{y}| \leq \norm{x}\norm{y}.
% \]
% \end{theorem}

\newtheorem{theorem}[definition]{Theorem}


% ------------------------------------------------------------
% PROPOSITION
% ------------------------------------------------------------
% Use for a useful mathematical result that is usually less
% central than a theorem.
%
% Usage:
%
% \begin{proposition}
% Every finite-dimensional subspace of \RR^n is closed.
% \end{proposition}

\newtheorem{proposition}[definition]{Proposition}


% ------------------------------------------------------------
% LEMMA
% ------------------------------------------------------------
% Use for an auxiliary result whose main purpose is to help
% prove another theorem/proposition.
%
% Usage:
%
% \begin{lemma}
% ...
% \end{lemma}

\newtheorem{lemma}[definition]{Lemma}




% ------------------------------------------------------------
% COROLLARY
% ------------------------------------------------------------
% Use for a result that follows relatively directly from a
% theorem/proposition that was just established.
%
% Usage:
%
% \begin{corollary}
% ...
% \end{corollary}

\newtheorem{corollary}[definition]{Corollary}


% ============================================================
% PROOF
% ============================================================
% The proof environment comes from amsthm.
%
% Usage:
%
% \begin{proof}
% Suppose that ...
%
% Therefore, ...
% \end{proof}
%
% LaTeX automatically adds the QED symbol at the end.


% ============================================================
% CUSTOM BOX: IDEA
% ============================================================
% Use this when you want to record intuition or the main idea
% behind a mathematical concept.
%
% Usage:
%
% \begin{idea}
% \textbf{Key idea.}
% A metric gives us a notion of distance...
% \end{idea}

\newmdenv[
    linewidth=1pt,
    skipabove=10pt,
    skipbelow=10pt
]{idea}


% ============================================================
% CUSTOM BOX: WARNING
% ============================================================
% Use this for common mistakes, distinctions, or things that
% are easy to confuse.
%
% Usage:
%
% \begin{warning}
% Do not confuse continuity with uniform continuity.
% \end{warning}

\newmdenv[
    linewidth=1pt,
    skipabove=10pt,
    skipbelow=10pt
]{warning}

% ============================================================
% CUSTOM BOX: NOTE
% ============================================================
% Use this for a short personal note, useful observation,
% connection between concepts, or explanation.
%
% Unlike "Remark", this is intended for informal notes.
%
% Usage:
%
% \begin{note}
% A sequence can be viewed as a function from $\NN$ to $\RR$.
% \end{note}

\newmdenv[
    linewidth=1pt,
    skipabove=10pt,
    skipbelow=10pt
]{note}


% ============================================================
% DOCUMENT
% ============================================================

\begin{document}


% ============================================================
% TITLE
% ============================================================

\begin{center}
    {\Large \textbf{Countable and Uncountable Sets}}\\[0.5em]
    {\large Lecture 7}
\end{center}

% Automatically generate the table of contents.
\tableofcontents

\newpage

\section{Countable and Uncountable}
\subsection{Recall: functions}
\begin{note}
This part, just recall injective, surjective, bijective of a \textbf{function} (not linear map, but it is the same anyway) in linear algebra. And if $A$ and $B$ are in 1-1 correspondence (exist a bijective function map from $A$ to $B$), we denote it as: $A \sim B$ (I think the lecture is abusing notation here).

\end{note}

\subsection{Counting}
\begin{definition}[Elementary Counting]
Counting, in fact, we are assigning natural number into elements of that set.

For example, we want to count oranges in a box, we simply point at each orange then say: "one, two, three,...". That's the assigning process.

So, mathematically, counting is mapping elements from $J_n=\{1,2,3,...,n\}$ (with $J_0 = \emptyset$) to the elements of the set that we want to count. This mapping is bijective.
\end{definition}

\begin{example}
Given $P =\{A, 1, :), =), @\}$, I can assign 1 to 1, $A$ to 3, :) to 2, @ to 4 and =) to 5. This is a valid assignment for this set. Thus, we have $P \sim J_5$. We can say: $|P| = 5$. And remember, we're mapping from $J_5$ to $P$, not the inverse (the lecture did it like that, but it is bijective mapping, I don't know if we can use the reverse way).
\end{example}

\begin{definition}[Finite]
We call $A$ finite if $A\sim J_n$, else $A$ is infinite.
\end{definition}

\begin{definition}[Countable]
We call $A$ countable if $A\sim \NN$.
\end{definition}

\begin{example}
Here are some examples of countable sets:
\begin{itemize}
    \item $\NN$ is countable. The bijective function $f: \NN \rightarrow \NN$ is $f(n) = n$.
    \item A sequence $x_1, x_2,x_3,...$ with distinct terms is countable. The bijective function $f: \NN \rightarrow seq$ is $f(n) = x_n$. This generalize to a set that can be listed in sequence is countable.
    \item $\{2,3,4,5,...\}$ is countable with $f(n) = n + 1$.
    \item $\{1,2,3,4,...,k-1,k+1,...\}$ is countable with $f(n) = n$ if $n<k$, $f(n) = n+1$ if $n\ge k$.
\end{itemize}
\end{example}

\begin{theorem}
$\NN$ is infinite.
\end{theorem}

\begin{note}
In lecture, they gonna proof the theorem by induction. I don't think I gonna mention it here. However, I gonna talk about the intuition (remember, the intuition is rejected by the lecturer, but I think it is still worth because of its logical imagination). If $\NN$ is finite, then we must have a bijective mapping from any $J_n$ to $\NN$. But no matter how you assign - or you can just assign like normal people (which is like counting) (assigning process is the "breaking point", lecturer says that there can be some assignments that we don't know that can make $\NN$ is 1-1 with some $J_n$, that's why they use induction?), you will always have number n+1 that has no "slot" (number) left to be assigned, that's why $\NN$ can't be 1-1 to any $J_n$.
\end{note}

\begin{example}
The interesting thing is, somehow, a set, can be 1-1 correspondence to its \textbf{proper} subset.

We have $\NN \sim 2\NN$, with $2\NN$ being the set of natural numbers. The bijective function is $f(n) = 2n$.

We say that, $\NN$ and $2\NN$ have the same cardinality ("size").
\end{example}

\begin{example}[$\ZZ$ is countable]
The bijective function is $f(0) = 1$, $f(2k) = k$, $f(2k+1) = -k$.
\end{example}

\begin{note}
The key thing in this whole counting scenario is the second example in Example 1.5. If a set can be listed in a way that you don't miss any element of it, then the set is countable or we have that "bijective function". In $\NN$ the listing is very easy, just $\{1,2,3,4,5,...\}$. In $\ZZ$, the bijective function above just do the following: 0 is assign to 1, 1 to 2, -1 to 3, 2 to 4, -2 to 5,... it go left right left right,... and assign numbers like that. In next section, you gonna see $\QQ_+$, which has an interesting way of listing elements of it, by using 2D array, 1-indexed, the $ij^{th}$ element is $\frac{i}{j}$, then you count zig-zag-ly, not ordered like $\NN$. Lastly, you can recall to the lecture of "Probabilistic Systems Analysis and Applied Probability", the lecturer said that if something is countable, you must have a way to list, then count it without missing a single element of the set.
\end{note}

\begin{theorem}
Any infinite subset $E$ of a countable set $A$ is countable.
\end{theorem}

\begin{theorem}
$\QQ$ is countable.    
\end{theorem}

\begin{note}
With the note above, you can convince yourself that $\QQ_+$ is indeed countable. Then we use the "left-right" trick of $\ZZ$ to prove the countability of $\QQ$.
\end{note}

\begin{theorem}
$\RR$ is uncountable.
\end{theorem}

\begin{note}
The proof of uncountable $\RR$ is assume $(0,1)$ is countable, then it have some kind of perfect assignment for $(0,1)$, but with some adjustments (Cantor), we gonna obtain a $y$ which in $(0,1)$ but not in the listing. As the rule above, the "prefect" assignment missed an element, hence $(0,1)$ is uncountable $\Rightarrow$ $\RR$ is uncountable.

Note that, if we apply the proof procedure of $\RR$ for $\QQ$, we can't proof $\QQ$ is uncountable. Given that any rational number, after the comma, they may have: finite number, infinite but in a pattern number. With the mix of that, if we use Cantor, we may obtain a infinite but no pattern after comma number, which is in $\RR$, not $\QQ$. Therefore the proof is broken.
\end{note}

\begin{theorem}
Given a set $A$, $\mathcal{P}(A)$ is the power set (set of all subsets of $A$). Then, $A \not\sim \mathcal{P}(A)$.
\end{theorem}

\begin{note}
The theorem above can be phrase like this:$|A| < |\mathcal{P}(A)|$. But why "$<$"? Given $A$, we have an injective map from $A \rightarrow \mathcal{P}(A): f(a) = \{a\} \in \mathcal{P}(A)$. Because there exists an injective map from $A \rightarrow \mathcal{P}(A)$, then $|A| \le |\mathcal{P}(A)|$. But there is no bijective map from $A \rightarrow \mathcal{P}(A)$, then $|A| < |\mathcal{P}(A)|$. The theorem tell us that, no matter how big is an infinite set, we can always create a bigger infinite set by taking its power set.
\end{note}

\end{document}